\begin{statementbox}
\textbf{C010\quad 椭圆切线距离公式（中心距、焦点积、主圆垂足）}

设椭圆
\[
C:\ \frac{x^2}{a^2}+\frac{y^2}{b^2}=1,\qquad a>b>0,
\qquad c=\sqrt{a^2-b^2},
\]
中心为 $O(0,0)$，焦点为 $F_1(-c,0)$、$F_2(c,0)$。$l$ 为椭圆的任意一条\textbf{切线}；$d_0,d_1,d_2$ 分别为 $O,F_1,F_2$ 到 $l$ 的距离，$H_1,H_2$ 分别为两焦点在 $l$ 上的垂足。

\begin{enumerate}
\item \textbf{中心到切线的距离（原心距）}\par
非竖直切线 $l:y=kx+m$ 必须满足
\[
m^2=a^2k^2+b^2,\qquad
\boxed{d_0=\frac{\sqrt{a^2k^2+b^2}}{\sqrt{1+k^2}}}.
\]
对竖直切线 $x=\pm a$，有 $d_0=a$。任意切线均满足 $\boxed{b\le d_0\le a}$。

\item \textbf{两焦点到同一切线的距离积}\par
\[
\boxed{d_1d_2=b^2},\qquad
\boxed{d_1+d_2=2d_0}.
\]
因此 $d_1,d_2$ 是方程 $t^2-2d_0t+b^2=0$ 的两根；特别地，$d_1=d_2=b$ 当且仅当切线为 $y=\pm b$。

\item \textbf{焦点垂足与主圆（长轴辅助圆）}\par
\[
\boxed{OH_1=OH_2=a},\qquad
\boxed{H_1,H_2\in\{(x,y):x^2+y^2=a^2\}}.
\]
固定其中任意一个焦点，让切线遍历椭圆的所有切线，其垂足的\textbf{完整轨迹}正是主圆 $x^2+y^2=a^2$，并非仅为主圆的一部分。

\item \textbf{通用直线形式（一次覆盖竖直切线）}\par
对 $l:Ax+By=D$，其中 $A^2+B^2>0$，它与 $C$ 相切的充要条件是
\[
\boxed{D^2=a^2A^2+b^2B^2}.
\]
此时 $d_0=\frac{|D|}{\sqrt{A^2+B^2}}$，$d_1=\frac{|D+Ac|}{\sqrt{A^2+B^2}}$，$d_2=\frac{|D-Ac|}{\sqrt{A^2+B^2}}$。
\end{enumerate}

\textbf{一句话记忆：}焦点距离\textbf{乘积看短轴}（$b^2$）；焦点垂足\textbf{落在长轴圆}（半径 $a$）；中心距离介于 $b$ 与 $a$ 之间。
\end{statementbox}
